\chapter{Технологический раздел}
В данном разделе описаны средства реализации и программная реализация
основных компонентов многопоточногосервера.

\section{Средства реализации}
Для реализации сервера был выбран язык программирования C стандарта C17.
Выбор обусловлен следующими причинами:
\begin{itemize}
	\item прямой доступ к системным вызовам Linux (epoll, sendfile);
	\item полный контроль над управлением памятью;
	\item совместимость с POSIX Threads.
\end{itemize}

Сборка проекта осуществляется с помощью утилиты GNU Make.
Компиляция производится компилятором GCC с флагами
\texttt{-Wall -Wextra -Wpedantic -std=c17}.
В режиме отладки используется флаг \texttt{-g} и макрос \texttt{DEBUG},
в режиме выпуска "--- флаг оптимизации \texttt{-O2}.

\section{Реализация алгоритма работы потока слушателя}
Поток слушатель выполняет приём входящих соединений
и регистрацию клиентских событий через подсистему epoll.
Каждый поток слушатель создаёт собственный серверный сокет
с опцией \texttt{SO\_REUSEPORT}, что позволяет ядру Linux
распределять входящие соединения между несколькими слушателями.

\subsection{Установка неблокирующего режима}
На листинге~\ref{lst:set_nonblocking} представлена функция перевода
файлового дескриптора в неблокирующий режим с помощью системного вызова
\texttt{fcntl}.

\begin{lstlisting}[caption={Установка неблокирующего режима},label={lst:set_nonblocking},language=C]
int setNonblocking(int fd)
{
    int flags = fcntl(fd, F_GETFL, 0);
    if (flags == -1)
        return 10;
    if (fcntl(fd, F_SETFL, flags | O_NONBLOCK) == -1)
        return 20;
    return 0;
}
\end{lstlisting}

\subsection{Обработка входящих подключений}
Функция \texttt{connectionHandle}, представленная на листинге~\ref{lst:connection_handle},
вызывается при готовности серверного сокета.


\begin{lstlisting}[caption={Обработка входящих подключений},label={lst:connection_handle},language=C]
void connectionHandle(int serverfd, int epollfd)
{
    while (1)
    {
        struct sockaddr_in clientAddr = {0};
        socklen_t clientLen = sizeof(clientAddr);
        int clientfd = accept(serverfd,
            (struct sockaddr *)&clientAddr, &clientLen);

        if (clientfd < 0)
        {
            if (errno == EAGAIN || errno == EWOULDBLOCK)
                break;
            break;
        }

        if (setNonblocking(clientfd) != 0)
        {
            close(clientfd);
            continue;
        }

        Task_t *task = createTask(clientfd, epollfd);
        if (!task)
        {
            close(clientfd);
            continue;
        }

        inet_ntop(AF_INET, &clientAddr.sin_addr,
                  task->clientip, sizeof(task->clientip));

        struct epoll_event ev = {0};
        ev.data.ptr = task;
        ev.events = EPOLLIN | EPOLLET
                  | EPOLLRDHUP | EPOLLONESHOT;

        if (epoll_ctl(epollfd, EPOLL_CTL_ADD,
                      clientfd, &ev) == -1)
        {
            deleteTask(task);
            close(clientfd);
            continue;
        }
    }
}
\end{lstlisting}

\subsection{Реализация алгоритма потока слушателя}
На листинге~\ref{lst:listener_thread} представлена функция потока слушателя.


\begin{lstlisting}[caption={Реализация потока слушателя},label={lst:listener_thread},language=C]
void *listenerThread(void *args)
{
    ServerConfig_t *conf = (ServerConfig_t *)args;

    int serverfd = socket(AF_INET, SOCK_STREAM, 0);
    if (serverfd < 0)
        return NULL;

    int opt = 1;
    setsockopt(serverfd, SOL_SOCKET, SO_REUSEPORT,
               &opt, sizeof(opt));

    struct sockaddr_in addr = {0};
    addr.sin_family = AF_INET;
    addr.sin_port = htons(conf->port);
    addr.sin_addr.s_addr = INADDR_ANY;

    bind(serverfd, (struct sockaddr *)&addr, sizeof(addr));
    listen(serverfd, MAXCON);
    setNonblocking(serverfd);

    int epollfd = epoll_create1(0);

    struct epoll_event ev = {0};
    ev.events = EPOLLIN | EPOLLET;
    ev.data.fd = serverfd;
    epoll_ctl(epollfd, EPOLL_CTL_ADD, serverfd, &ev);

    /* Таймер статистики (каждые 5 секунд) */
    int timerfd = timerfd_create(CLOCK_MONOTONIC,
                                 TFD_NONBLOCK);
    struct itimerspec timer_spec = {0};
    timer_spec.it_value.tv_sec = 5;
    timer_spec.it_interval.tv_sec = 5;
    timerfd_settime(timerfd, 0, &timer_spec, NULL);
    ev.data.fd = timerfd;
    ev.events = EPOLLIN;
    epoll_ctl(epollfd, EPOLL_CTL_ADD, timerfd, &ev);

    /* eventfd для сигнала завершения */
    ev.data.fd = conf->shutdownfd;
    ev.events = EPOLLIN;
    epoll_ctl(epollfd, EPOLL_CTL_ADD,
              conf->shutdownfd, &ev);

    struct epoll_event events[MAXEVENTS];

    while (!conf->shutdown)
    {
        int nfds = epoll_wait(epollfd, events,
                              MAXEVENTS, 1000);
        if (nfds == -1)
        {
            if (errno == EINTR) continue;
            break;
        }

        for (int i = 0; i < nfds; ++i)
        {
            if (events[i].data.fd == conf->shutdownfd)
                goto cleanup;
            else if (events[i].data.fd == serverfd)
                connectionHandle(serverfd, epollfd);
            else if (events[i].data.fd == timerfd)
                timerHandle(timerfd);
            else
            {
                Task_t *task =
                    (Task_t *)events[i].data.ptr;
                if (events[i].events
                    & (EPOLLRDHUP | EPOLLHUP))
                {
                    epoll_ctl(epollfd, EPOLL_CTL_DEL,
                              task->clientfd, NULL);
                    close(task->clientfd);
                    deleteTask(task);
                    continue;
                }
                push(conf->queue, task);
            }
        }
    }

cleanup:
    close(timerfd);
    close(epollfd);
    close(serverfd);
    return NULL;
}
\end{lstlisting}

\section{Реализация алгоритма работы потока обработчика}
Поток-обработчик извлекает задачи из очереди и обрабатывает HTTP-запросы
в соответствии с конечным автоматом клиентского соединения.

\begin{lstlisting}[caption={Реализация функции чтения HTTP-запроса},label={lst:reading_handle},language=C]
int readingHandle(Task_t *task)
{
    while (1)
    {
        if (task->readPos >= task->readLen - 1)
            return -1;

        ssize_t n = recv(task->clientfd,
                        task->readBuf + task->readPos,
                        task->readLen - task->readPos - 1,
                        MSG_DONTWAIT);
        if (n > 0)
        {
            task->readPos += n;
            task->readBuf[task->readPos] = '\0';
            task->lastActivity = time(NULL);
            if (strstr(task->readBuf, "\r\n\r\n"))
            {
                task->state = SENDING_HEADERS;
                return 1;
            }
        }
        else if (n == 0)
            return -1;
        else
        {
            if (errno == EAGAIN || errno == EWOULDBLOCK)
                return 0;
            return -1;
        }
    }
}
\end{lstlisting}

\subsection{Реализация функции парсинга HTTP-запроса}
На листинге~\ref{lst:parse_request} представлена функция разбора HTTP-запроса.

\begin{lstlisting}[caption={Реализация функции парсинг HTTP-запроса},label={lst:parse_request},language=C]
int parseRequest(Task_t *task)
{
    char *buf = task->readBuf;

    if (strncmp(buf, "GET ", 4) == 0)
    {
        task->method = GET;
        buf += 4;
    }
    else if (strncmp(buf, "HEAD ", 5) == 0)
    {
        task->method = HEAD;
        buf += 5;
    }
    else
    {
        task->method = UNKNOWN;
        task->statusCode = 405;
        return 0;
    }

    char *pathStart = buf;
    char *pathEnd = strchr(buf, ' ');
    if (!pathEnd)
    {
        task->statusCode = 400;
        return -1;
    }

    size_t pathLen = pathEnd - pathStart;
    task->targetPath = malloc(pathLen + 1);
    if (!task->targetPath)
    {
        task->statusCode = 500;
        return -1;
    }
    memcpy(task->targetPath, pathStart, pathLen);
    task->targetPath[pathLen] = '\0';

    if (strstr(task->targetPath, ".."))
    {
        task->statusCode = 403;
        return 0;
    }

    char *connHeader =
        strcasestr(task->readBuf, "\r\nConnection:");
    if (connHeader)
    {
        connHeader += 13;
        while (*connHeader == ' ') connHeader++;
        if (strncasecmp(connHeader, "close", 5) == 0)
            task->keepAlive = 0;
    }
    return 0;
}
\end{lstlisting}

\subsection{Реализация функции формирования HTTP-ответа}
На листинге~\ref{lst:prepare_response} представлена функция формирования
HTTP-ответа. 
\begin{lstlisting}[caption={Реализация функции формирования HTTP-ответа},label={lst:prepare_response},language=C]
int prepareResponse(Task_t *task)
{
    char filepath[512];
    struct stat st;
    const char *mimeType = "text/html";
    const char *body = NULL;
    size_t bodyLen = 0;

    if (task->statusCode == 200)
    {
        const char *target =
            (strcmp(task->targetPath, "/") == 0)
            ? "/index.html" : task->targetPath;
        snprintf(filepath, sizeof(filepath),
                 "%s%s", WWW_ROOT, target);

        task->filefd = open(filepath, O_RDONLY);
        if (task->filefd < 0)
            task->statusCode = 404;
        else if (fstat(task->filefd, &st) < 0
                 || !S_ISREG(st.st_mode))
        {
            close(task->filefd);
            task->filefd = -1;
            task->statusCode = 404;
        }
        else
        {
            task->fileSize = st.st_size;
            mimeType = getMimeType(filepath);
        }
    }

    if (task->statusCode != 200)
    {
        if (task->filefd >= 0)
        {
            close(task->filefd);
            task->filefd = -1;
        }
        task->fileSize = 0;
        switch (task->statusCode)
        {
            case 403: body = ERROR_403_BODY; break;
            case 404: body = ERROR_404_BODY; break;
            case 405: body = ERROR_405_BODY; break;
            default:  body = ERROR_500_BODY; break;
        }
        bodyLen = strlen(body);
    }

    size_t contentLength =
        (task->filefd >= 0) ? task->fileSize : bodyLen;

    char headers[1024];
    int headersLen = snprintf(headers, sizeof(headers),
        "HTTP/1.1 %d %s\r\n"
        "Content-Type: %s\r\n"
        "Content-Length: %zu\r\n"
        "Connection: %s\r\n\r\n",
        task->statusCode,
        getStatusText(task->statusCode),
        mimeType, contentLength,
        task->keepAlive ? "keep-alive" : "close");

    size_t totalSize = headersLen
                     + (body ? bodyLen : 0);
    task->writeBuf = malloc(totalSize);
    if (!task->writeBuf)
        return -1;

    memcpy(task->writeBuf, headers, headersLen);
    if (body)
        memcpy(task->writeBuf + headersLen,
               body, bodyLen);

    task->writeLen = totalSize;
    task->writePos = 0;
    return 0;
}
\end{lstlisting}

\subsection{Реализация функции отправки HTTP-заголовков}
На листинге~\ref{lst:sending_headers} представлена функция отправки
HTTP-заголовков клиенту.

\begin{lstlisting}[caption={Реализация функции отправки HTTP-заголовков},label={lst:sending_headers},language=C]
int sendingHeadersHandle(Task_t *task)
{
    while (task->writePos < task->writeLen)
    {
        ssize_t n = send(task->clientfd,
                         task->writeBuf + task->writePos,
                         task->writeLen - task->writePos,
                         MSG_DONTWAIT);
        if (n > 0)
            task->writePos += n;
        else if (n < 0)
        {
            if (errno == EAGAIN || errno == EWOULDBLOCK)
                return 0;
            return -1;
        }
    }

    if (task->filefd >= 0 && task->method == GET)
        task->state = SENDING_BODY;
    else
        task->state = DONE;
    return 1;
}
\end{lstlisting}

\subsection{Реализация функции отправки тела ответа}
На листинге~\ref{lst:sending_body} представлена функция отправки
тела ответа. 
\begin{lstlisting}[caption={Реализация функции отправки тела ответа через sendfile},label={lst:sending_body},language=C]
int sendingBodyHandle(Task_t *task)
{
    while (task->filePos < task->fileSize)
    {
        off_t offset = task->filePos;
        ssize_t sent = sendfile(task->clientfd,
            task->filefd, &offset,
            task->fileSize - task->filePos);

        if (sent > 0)
            task->filePos += sent;
        else if (sent < 0)
        {
            if (errno == EAGAIN
                || errno == EWOULDBLOCK)
                return 0;
            return -1;
        }
    }
    task->state = DONE;
    return 1;
}
\end{lstlisting}

\subsection{Реализация функции перерегистрации дескриптора в epoll}
На листинге~\ref{lst:rearm_epoll} представлена функция перерегистрации
дескриптора в epoll после обработки. 
\begin{lstlisting}[caption={Реализация функции перерегистрации дескриптора в epoll},label={lst:rearm_epoll},language=C]
void rearmEpoll(Task_t *task, uint32_t events)
{
    struct epoll_event ev = {0};
    ev.data.ptr = task;
    ev.events = events | EPOLLET
              | EPOLLRDHUP | EPOLLONESHOT;
    epoll_ctl(task->epollfd, EPOLL_CTL_MOD,
              task->clientfd, &ev);
}
\end{lstlisting}

\subsection{Реализация функции потока обработчика}
На листинге~\ref{lst:worker_thread} представлена функция потока обработчика.
Поток в цикле извлекает задачи из очереди и обрабатывает их
в соответствии с текущим состоянием конечного автомата.
Переходы между состояниями:
\texttt{READING} $\rightarrow$ \texttt{SENDING\_HEADERS} $\rightarrow$ \texttt{SENDING\_BODY} $\rightarrow$ \texttt{DONE}.
При поддержке keep-alive после состояния \texttt{DONE}
соединение возвращается в состояние \texttt{READING}
для обработки следующего запроса.
При ошибке \texttt{EAGAIN} дескриптор перерегистрируется в epoll,
и задача возвращается в пул событий.

\begin{lstlisting}[caption={Реализация функции потока обработчика},label={lst:worker_thread},language=C]
void *workerThread(void *args)
{
    ServerConfig_t *conf = (ServerConfig_t *)args;
    int needClose;

    while (!conf->shutdown)
    {
        Task_t *task = NULL;
        needClose = 0;

        if (pop(conf->queue, &task) != 0)
            break;
        if (!task)
            continue;

        atomic_fetch_add(&total_requests, 1);

        switch (task->state)
        {
        case READING:
        {
            int rc = readingHandle(task);
            if (rc == 1)
            {
                if (parseRequest(task) < 0
                    || prepareResponse(task) < 0)
                {
                    needClose = 1; break;
                }
                rc = sendingHeadersHandle(task);
                if (rc == 0)
                {
                    rearmEpoll(task, EPOLLOUT);
                    continue;
                }
                else if (rc < 0)
                {
                    needClose = 1; break;
                }
                if (task->state == SENDING_BODY)
                {
                    rc = sendingBodyHandle(task);
                    if (rc == 0)
                    {
                        rearmEpoll(task, EPOLLOUT);
                        continue;
                    }
                    else if (rc < 0)
                    {
                        needClose = 1; break;
                    }
                }
            }
            else if (rc == 0)
            {
                rearmEpoll(task, EPOLLIN);
                continue;
            }
            else
            {
                needClose = 1; break;
            }
        }
        /* fallthrough */
        case SENDING_HEADERS:
        {
            int rc = sendingHeadersHandle(task);
            if (rc == 0)
            {
                rearmEpoll(task, EPOLLOUT);
                continue;
            }
            else if (rc < 0)
            {
                needClose = 1; break;
            }
            if (task->state != SENDING_BODY)
                break;
        }
        /* fallthrough */
        case SENDING_BODY:
        {
            int rc = sendingBodyHandle(task);
            if (rc == 0)
            {
                rearmEpoll(task, EPOLLOUT);
                continue;
            }
            else if (rc < 0)
            {
                needClose = 1; break;
            }
        }
        /* fallthrough */
        case DONE:
            if (task->keepAlive)
            {
                close(task->filefd);
                task->filefd = -1;
                free(task->writeBuf);
                task->writeBuf = NULL;
                free(task->targetPath);
                task->targetPath = NULL;
                task->state = READING;
                task->readPos = 0;
                task->writeLen = 0;
                task->writePos = 0;
                task->fileSize = 0;
                task->filePos = 0;
                task->statusCode = 200;
                task->method = UNKNOWN;
                rearmEpoll(task, EPOLLIN);
                continue;
            }
            needClose = 1;
            break;
        }

        if (needClose)
        {
            epoll_ctl(task->epollfd, EPOLL_CTL_DEL,
                      task->clientfd, NULL);
            close(task->clientfd);
            deleteTask(task);
        }
    }
    return NULL;
}
\end{lstlisting}

\section{Реализация алгоритма работы главного потока сервера}
Главный поток сервера выполняет инициализацию всех подсистем
и ожидает сигнал завершения. На листинге~\ref{lst:main_server}
представлена функция \texttt{main}

\begin{lstlisting}[caption={Реализация функции главного потока сервера},label={lst:main_server},language=C]
int main(int argc, char **argv)
{
    ServerConfig_t conf = {0};
    int rc = parseArgs(argc, argv, &conf);
    if (rc != 0)
        return (rc == 1) ? EXIT_SUCCESS : EXIT_FAILURE;

    signal(SIGPIPE, SIG_IGN);
    init_logger("logs/server.log");

    sigset_t sigset;
    sigemptyset(&sigset);
    sigaddset(&sigset, SIGINT);
    sigaddset(&sigset, SIGTERM);
    pthread_sigmask(SIG_BLOCK, &sigset, NULL);

    conf.queue = createQueue();
    conf.shutdown = 0;
    conf.shutdownfd = eventfd(0, EFD_NONBLOCK);

    pthread_t *workers =
        calloc(conf.workers, sizeof(pthread_t));
    pthread_t *listeners =
        calloc(conf.listeners, sizeof(pthread_t));

    /* Создание пулов потоков */
    for (int i = 0; i < conf.workers; ++i)
        pthread_create(&workers[i], NULL,
                       &workerThread, &conf);
    for (int i = 0; i < conf.listeners; ++i)
        pthread_create(&listeners[i], NULL,
                       &listenerThread, &conf);

    /* Ожидание сигнала завершения */
    int sig;
    sigwait(&sigset, &sig);

    /* Корректное завершение */
    conf.shutdown = 1;
    shutdownQueue(conf.queue);

    uint64_t val = 1;
    write(conf.shutdownfd, &val, sizeof(val));

    for (int i = 0; i < conf.listeners; ++i)
        pthread_join(listeners[i], NULL);
    for (int i = 0; i < conf.workers; ++i)
        pthread_join(workers[i], NULL);

    free(workers);
    free(listeners);
    close(conf.shutdownfd);
    deleteQueue(conf.queue);
    close_logger();
    return EXIT_SUCCESS;
}
\end{lstlisting}

